\section{Discussion}
\label{sec:discussion}

The central empirical finding of this study is that reliability, not raw model correctness, sustains a working large-scale agentic simulation. Retained model autonomy was 0.00--0.02\% across every system, workload, and configuration we evaluated, while 99.99\% of steps remained useful. A real 7--8B instruct model produced a raw output that failed validation, even after an in-loop repair attempt, roughly half the time. It is the pipeline's policy-completion layer, not the model, that keeps the simulation moving. We read this as a demonstrated property of the model class and workloads evaluated here, not a general claim about all instruction-tuned models: a substantially larger or differently prompted model could plausibly raise the retained-autonomy rate, and testing that directly is future work.

A second finding concerns measurement methodology itself. Three independent, structurally identical bugs each hid or inverted a real effect in this study: an evaluation harness that never generated enough real concurrent load to stress a serving parameter under test, a hardcoded per-round batching cap that added a second, independent concurrency ceiling, and an in-flight concurrency limit tuned under the same insufficient-load regime as the first bug. We found each through the same diagnostic move, tracing the call path to the parameter that actually bound execution rather than trusting the parameter that was intended to. We interpret this as a transferable risk for anyone benchmarking LLM-serving or agent-scheduling configurations under real infrastructure: a statistically well-powered, non-overlapping result is not evidence of a real effect unless the load level has been confirmed to stress the parameter under test. This is a supported interpretation drawn from a single investigation arc, not a claim we have tested across other codebases or serving stacks.

The scheduling decision gate's failure to clear is the case where the demonstrated result and its interpretation must be kept most clearly separate. What we demonstrated is that all six concurrent scheduling strategies are statistically indistinguishable on both systems once the three measurement confounds are corrected. What we do not conclude is that capability-aware or queue-aware scheduling is ineffective. Both mechanisms are keyed on a per-request provider label, and every real-hardware run in this study routed every agent to the same single serving endpoint, so the workloads evaluated cannot exercise the property either mechanism is designed to exploit. Whether the gate would clear under genuine multi-provider heterogeneity, for instance two concurrently live servers with different real capacity limits, remains an open, testable question, and building that evaluation is the most direct next step this study identifies.

The one cross-system reliability difference we observed, a consistently lower invalid-output rate on the NVIDIA system than the AMD system for the same workload, model revision, and cached-attention precision, illustrates the same discipline applied to a smaller finding. The two systems also differ by one minor serving-software release and in the attention-kernel implementation actually exercised. Either difference could plausibly explain a generation-quality gap of this size, so we report the observation without attributing it to the accelerator vendor.
